\begin{examplebox}
\par\noindent\begin{minipage}{\linewidth}
\textbf{例1：最简单的坐标题——不求角平分线方程}

\textbf{题目：} $A(0,3),B(4,0),C(0,0)$，求 $\triangle ABC$ 的内心 $I$。
\VIncenterExample{1}

\end{minipage}\par
\textbf{解：}
\begin{enumerate}
\item 先算\textbf{对边}：$a=BC=4$，$b=CA=3$，$c=AB=\sqrt{4^2+3^2}=5$。
\item 分母是周长 $a+b+c=12$。逐项代入：
\[
\begin{aligned}
x_I&=\frac{4\cdot0+3\cdot4+5\cdot0}{12}=1,\\
y_I&=\frac{4\cdot3+3\cdot0+5\cdot0}{12}=1.
\end{aligned}
\]
\item \textbf{检验：} 这是两条直角边分别在坐标轴上的直角三角形；内切圆圆心应形如 $(r,r)$。由面积 $S=3\cdot4/2=6$ 和半周长 $p=6$，得 $r=S/p=1$，与 $(1,1)$ 一致。
\end{enumerate}
\textbf{答案：}$\boxed{I(1,1)}$。\textbf{省掉的工作：}不用分别求两条角平分线的方程再联立。
\end{examplebox}

\begin{examplebox}
\par\noindent\begin{minipage}{\linewidth}
\textbf{例2：反向识别——系数是边长的倍数也可以}

\textbf{题目：}已知 $AB=3,AC=4,BC=5$，点 $P$ 满足
\[
10\vec{PA}+8\vec{PB}+6\vec{PC}=\vec0.
\]
判断点 $P$。
\VIncenterExample{2}

\end{minipage}\par
\textbf{解：}
\begin{enumerate}
\item 检查三个\textbf{对边}：$a=5,b=4,c=3$。
\item 题目中的系数 $10:8:6=5:4:3=a:b:c$。将题目等式两边除以 $2$，得到
\[5\vec{PA}+4\vec{PB}+3\vec{PC}=\vec0.\]
\item 内心 $I$ 满足同一个等式。将任意参考点 $O$ 引入，解出
\[
\vec{OP}=\frac{5\vec{OA}+4\vec{OB}+3\vec{OC}}{12}
=\vec{OI}.
\]
同一参考点到两点的向量相等，说明 $P=I$。
\end{enumerate}
\textbf{答案：}$\boxed{P\text{ 是内心}}$。\textbf{关键：}必须用对边配对，不能仅凭“三个系数不相等”就判断。
\end{examplebox}

\begin{examplebox}
\par\noindent\begin{minipage}{\linewidth}
\textbf{例3：只求角平分线与边的交点——不用先求内心}

\textbf{题目：} $A(0,0),B(4,0),C(0,3)$，$AD$ 为 $\angle A$ 的角平分线且 $D\in BC$，求 $D$ 及 $\vec{AD}$。
\VIncenterExample{3}

\end{minipage}\par
\textbf{解：}
\begin{enumerate}
\item $AB=c=4$，$AC=b=3$。由角平分线定理 $BD:DC=4:3$，所以 $BD/BC=4/7$。
\item 从 $B$ 沿 $\vec{BC}=(-4,3)$ 走 $4/7$，可得
\[
D=B+\frac47(C-B)=(4,0)+\frac47(-4,3)=\left(\frac{12}{7},\frac{12}{7}\right).
\]
\item 因为 $A=(0,0)$，$\vec{AD}=(12/7,12/7)$。也可直接用
\[
\vec{AD}=\frac37(4,0)+\frac47(0,3)=\left(\frac{12}{7},\frac{12}{7}\right).
\]
\end{enumerate}
\textbf{答案：}$\boxed{D(12/7,12/7)}$。
\textbf{解题选择：}只需要角平分线交点 $D$，就用角平分线定理；不要绕路先求内心。
\end{examplebox}

\begin{examplebox}
\par\noindent\begin{minipage}{\linewidth}
\textbf{例4：角平分线方向——单位向量相加有什么用？}

\textbf{题目：} 在例3的三角形中，直接写出从 $A$ 出发的内角平分线方向，并说明它是否能单独确定内心坐标。
\VIncenterExample{4}

\end{minipage}\par
\textbf{解：}
\begin{enumerate}
\item $\vec{AB}=(4,0)$，$|\vec{AB}|=4$，所以 $\vec{AB}/|\vec{AB}|=(1,0)$。
\item $\vec{AC}=(0,3)$，$|\vec{AC}|=3$，所以 $\vec{AC}/|\vec{AC}|=(0,1)$。
\item 相加得到 $(1,1)$，因此从 $A$ 出发的内角平分线就是方向为 $(1,1)$ 的直线，即 $y=x$。
\item 但 $y=x$ 上有无数个点，\textbf{不能单独确定内心}。再求三边 $a=5,b=3,c=4$，代入内心坐标式：
\[
I=\frac{5(0,0)+3(4,0)+4(0,3)}{12}=(1,1).
\]
\end{enumerate}
\textbf{结论：}单位向量之和用来找\textbf{方向}；对边长乘顶点坐标才可直接找\textbf{位置}。
\end{examplebox}

\begin{examplebox}
\par\noindent\begin{minipage}{\linewidth}
\textbf{例5：同一道题，内心与重心不是同一点}

\textbf{题目：} $A(0,0), B(6,0), C(0,8)$，比较内心 $I$ 与重心 $G$ 的坐标。
\VIncenterExample{5}

\end{minipage}\par
\textbf{解：}
\begin{enumerate}
\item $a=BC=10,\ b=CA=8,\ c=AB=6$。
\item 内心按\textbf{对边长}计算：
\[
I=\frac{10(0,0)+8(6,0)+6(0,8)}{24}=(2,2).
\]
\item 重心对三个顶点\textbf{一视同仁}：
\[
G=\frac{(0,0)+(6,0)+(0,8)}{3}=\left(2,\frac83\right).
\]
\end{enumerate}
\textbf{比较：}两者横坐标恰好相同，纵坐标不同。不能把内心坐标误写成三个顶点坐标的算术平均。
\end{examplebox}

\begin{examplebox}
\par\noindent\begin{minipage}{\linewidth}
\textbf{例6：面积比可以先于向量使用}

\textbf{题目：}三角形三边 $BC=13,CA=14,AB=15$，内心为 $I$。求三个小三角形面积之比，并在已知原面积为 $84$ 时分别求其面积。
\VIncenterFigure{5}{.47}{13}{14}{15}{图内数字是三块面积，边上的数字是边长；同高 $r=4$，所以面积分别为 $26,28,30$。}

\end{minipage}\par
\textbf{解：}
\begin{enumerate}
\item 三块三角形从内心到对应边的高都相同，故
\[
S_{IBC}:S_{ICA}:S_{IAB}=13:14:15.
\]
\item 三个比数之和为 $42$，原面积为 $84$，每一份为 $84/42=2$。
\item 因此
\[
\boxed{S_{IBC}=26,\quad S_{ICA}=28,\quad S_{IAB}=30.}
\]
\item 验算：$26+28+30=84$。还可得内切圆半径 $r=2S/(a+b+c)=168/42=4$，例如 $S_{IBC}=13\cdot4/2=26$。
\end{enumerate}
\textbf{解题启示：}如果题目只问面积比，连向量都不需要引入。
\end{examplebox}
